\subsection{Attacker model}
\label{sec:tm-model}

We assume an attacker with:

\begin{itemize}
  \item the ability to open one or more TCP connections to the
    target server (no authentication, no shared key),
  \item the ability to send well-formed HTTP/1.1 requests with
    arbitrary chunked-TE bodies,
  \item modest egress bandwidth (KB/s, not MB/s; a single
    residential broadband connection suffices),
  \item no co-location requirement, no on-path position, no TLS
    private key.
\end{itemize}

The attacker does not need to bypass authentication. Per the
attack model (Section~\ref{sec:tm-reach}), the target endpoint is
\emph{any} HTTP path that accepts a request body; for most
deployments this includes unauthenticated paths.

\subsection{Reachability}
\label{sec:tm-reach}

Per-chunk amplification is reachable on any HTTP/1.1 endpoint that
accepts a request body, irrespective of method. POST, PUT, PATCH,
and DELETE-with-body all trigger the parser. The endpoint does
\emph{not} have to do anything with the body; even a handler that
returns 401 / 403 immediately on auth failure has already paid the
per-chunk cost by the time it inspects the body.

This makes the attack reachable on:

\begin{itemize}
  \item unauthenticated APIs (chat-as-a-service, contact forms,
    upload endpoints, webhook receivers, OAuth callback URLs),
  \item authenticated APIs that begin reading the body before the
    auth check completes (the default order in many ASGI / WSGI /
    Servlet stacks),
  \item any endpoint reached through a frontend that does not
    aggregate (Section~\ref{sec:mit-frontend}: HAProxy, some CDN
    configurations).
\end{itemize}

The wire pattern is RFC-compliant; intrusion detection systems
keyed on protocol violations will not fire.

\subsection{Cost-per-attacker-byte}
\label{sec:tm-cost}

Each one-byte chunk on the wire is 6 bytes
(\texttt{1$\backslash$r$\backslash$nA$\backslash$r$\backslash$n}).
For $N{=}\num{250000}$ chunks the attacker sends roughly 1.5 MB
of wire bytes plus headers; the server costs are summarized
in Table~\ref{tab:headline-costs}.

\begin{table}[h]
  \centering
  \footnotesize
  \caption{Per-request attacker bandwidth versus server cost,
  $N{=}\num{250000}$ one-byte chunks. The cost column mixes two
  measurement types, marked accordingly: \textsuperscript{c}~cells
  are Mode~B server-CPU-seconds
  ($T_\mathrm{wall}\times\texttt{cpu\%}$ from a \fname{docker stats}
  snapshot); \textsuperscript{w}~cells are single-core \emph{wall}
  seconds (server pinned at one vCPU, so wall is an upper bound on
  CPU; \textsuperscript{A}~denotes a Mode~A cell,
  \textsuperscript{e}~an extrapolated cell). Asymmetry (cost-seconds
  per attacker-MB-on-wire, higher favors the attacker) is exact only
  for \textsuperscript{c}~rows; see Section~\ref{sec:method-metric}.}
  \label{tab:headline-costs}
  \begin{tabularx}{\linewidth}{@{}l r r r@{}}
    \toprule
    Server & Wire (MB) & Cost (s) & Asymmetry (s/MB) \\
    \midrule
    Vert.x                               & 1.5 & 1.8\textsuperscript{c}   & 1.2 \\
    Go net/http                          & 1.5 & 3.4\textsuperscript{c}   & 2.3 \\
    Spring Boot / Tomcat                 & 1.5 & 5.3\textsuperscript{c}   & 3.5 \\
    nginx as origin                      & 1.5 & 29\textsuperscript{c}    & 19  \\
    bandit (linear)                      & 1.5 & 35\textsuperscript{c}    & 23  \\
    Kestrel 9                            & 1.5 & 25.9\textsuperscript{w}  & --- \\
    node \fwid{http.Server} (QUADRATIC)  & 1.5 & 282\textsuperscript{w,A} & --- \\
    express (QUADRATIC)                  & 1.5 & 324\textsuperscript{w,A,e} & --- \\
    \bottomrule
  \end{tabularx}

  {\footnotesize\textsuperscript{c}~Mode~B server CPU-seconds
  (\fname{docker stats} snapshot $\times$ wall).
  \textsuperscript{w}~wall-derived upper bound, not CPU.
  \textsuperscript{A}~Mode~A cell.
  \textsuperscript{e}~extrapolated from a quadratic fit, not
  measured. uvicorn rows are omitted: they have no Mode~B
  $N{=}\num{250000}$ CPU cell.}
\end{table}

A residential broadband uplink at \SI{20}{\mebi\bit\per\second}
sustains roughly \SI{2.5}{\mebi\byte\per\second}, comfortably
servicing one or two attacker connections per second. Against
\fwid{express} (\(N{=}\num{250000}\), Mode A, \emph{extrapolated}
cell: 324 s of single-core \emph{wall} time per request, not a
measured CPU integral), one such uplink sustained at one request
per second per source IP would, if the server stays CPU-bound on
one vCPU, draw on the order of a few hundred core-seconds per
wall-second from the target. We flag this as a wall-time-based
estimate built on an extrapolated cell: the directly-measured
Mode~B CPU cost of a single Node \fwid{http.Server} request at
$N{=}\num{250000}$ is far smaller (about \SI{1.5}{\second}), and
the gap between that and the saturated-core wall figure is exactly
why the "saturating any single-region deployment" claim requires a
measured-CPU derivation we leave for future work
(Section~\ref{sec:threats}, Section~\ref{sec:future-work}).

\subsection{Detection and evasion}
\label{sec:tm-detect}

\paragraph{Detection.}
The attack signature on the wire is "request with
\fname{Transfer-Encoding: chunked} and many small chunks." Most
WAFs do not include a rule for this; most observability stacks log
HTTP requests at the application layer (where the chunked-TE detail
is lost) rather than at the L4 layer (where per-packet timing
would expose it). The attack is functionally invisible to a
default observability stack.

A per-connection CPU-accounting layer (cgroup CPU controller, BPF
per-cgroup timing) would expose the bug, but is not commonly
deployed.

\paragraph{Evasion.}
A capable attacker can vary chunk size between 1 and a few bytes
(keeping the per-chunk cost dominant but breaking signature-based
detection on minimum chunk size). The attacker can rotate source
IPs to dodge per-IP rate limits, and can interleave the malicious
requests with legitimate-looking traffic on the same connection
under HTTP keep-alive.

\subsection{Combined attacks}
\label{sec:tm-combined}

The per-chunk amplification combines unfavorably with several
adjacent vulnerabilities:

\begin{itemize}
  \item \emph{slow-loris on headers} (CWE-400 class): an attacker
    can spend
    seconds sending request headers slowly before chunked-TE body
    parsing even begins, holding the connection-handling worker
    in head-of-line block.
  \item \emph{repeated-header coalescing} (Tornado
    CVE-2025-67725, Django ASGI CVE-2025-14550): if the target
    is also vulnerable to one of those, the same connection
    can trigger both costs.
  \item \emph{HTTP/2 multiplexing}: a single connection can carry
    many concurrent streams, each carrying chunked-TE bodies, all
    feeding the same per-chunk parser. We have not measured this
    interaction (Section~\ref{sec:future}) but it likely amplifies
    the per-connection cost by the maximum concurrent stream
    count.
\end{itemize}

\subsection{Distributed amplification}
\label{sec:tm-distributed}

A botnet scales the attack linearly with attacker count. We
caution that the per-request Node figures above are single-core
\emph{wall} time (and, for \fwid{express}, extrapolated), not
integrated CPU-seconds, so the following is an order-of-magnitude
estimate rather than a measured result. At one request per second
per IP against an \fwid{express}-like target, a 100-IP botnet
(trivial cost on commercial residential proxy services) would, if
the target stays CPU-bound on the affected workers, draw on the
order of $100\times$ the per-request core-time -- hundreds of
core-seconds per wall-second. The measured Mode~B CPU cost of a
single Node \fwid{http.Server} request at $N{=}\num{250000}$ is
about \SI{1.5}{\second}, roughly $200\times$ smaller than the
saturated-core wall figure; until a quadratic Node server is
re-measured with integrated cgroup CPU at this $N$
(Section~\ref{sec:threats}), we state the botnet draw only as an
upper-bound estimate and do not claim it categorically saturates a
multi-region deployment.
